%%%%%%%%%%%%%%%%
%%% deut 9 %%%
%%%%%%%%%%%%%%%%

\Chapter{9}

\Verse{1}
{
	\hDtr1{שְׁמַע יִשְׂרָאֵל אַתָּה עֹבֵר הַיּוֹם אֶת הַיַּרְדֵּן לָבֹא לָרֶשֶׁת גּוֹיִם גְּדֹלִים וַעֲצֻמִים מִמֶּךָּ עָרִים גְּדֹלֹת וּבְצֻרֹת בַּשָּׁמָיִם׃}
}
{
	\eDtr1{
		9 “Listen, Israel: you’re crossing the Jordan today to
		come to dispossess nations bigger and stronger than you,
		big cities and fortified to the skies,
	}
}
{
}

\Verse{2}
{
	\hDtr1{עַם גָּדוֹל וָרָם בְּנֵי עֲנָקִים אֲשֶׁר אַתָּה יָדַעְתָּ וְאַתָּה שָׁמַעְתָּ מִי יִתְיַצֵּב לִפְנֵי בְּנֵי עֲנָק׃}
}
{
	\eDtr1{
		a people big and tall, giants, of whom you have known, and
		of whom you have heard it said: ‘Who can stand in front of
		the giants?!’
	}
}
{
}

\Verse{3}
{
	\hDtr1{וְיָדַעְתָּ הַיּוֹם כִּי יְהֹוָה אֱלֹהֶיךָ הוּא הָעֹבֵר לְפָנֶיךָ אֵשׁ אֹכְלָה הוּא יַשְׁמִידֵם וְהוּא יַכְנִיעֵם לְפָנֶיךָ וְהוֹרַשְׁתָּם וְהַאֲבַדְתָּם מַהֵר כַּאֲשֶׁר דִּבֶּר יְהֹוָה לָךְ׃}
}
{
	\eDtr1{
		So you shall know today that YHWH, your God: He is the one
		who is crossing in front of you, a consuming fire. He will
		destroy them, and He will subdue them in front of you, so
		you’ll dispossess them and destroy them quickly as YHWH
		has spoken to you.
	}
}
{
}

\Verse{4}
{
	\hDtr1{אַל תֹּאמַר בִּלְבָבְךָ בַּהֲדֹף יְהֹוָה אֱלֹהֶיךָ אֹתָם מִלְּפָנֶיךָ לֵאמֹר בְּצִדְקָתִי הֱבִיאַנִי יְהֹוָה לָרֶשֶׁת אֶת הָאָרֶץ הַזֹּאת וּבְרִשְׁעַת הַגּוֹיִם הָאֵלֶּה יְהֹוָה מוֹרִישָׁם מִפָּנֶיךָ׃}
}
{
	\eDtr1{
		“Don’t say in your heart when YHWH pushes them from in
		front of you, saying, ‘Because of my virtue YHWH has
		brought me to take possession of this land,’ when it is
		because of these nations’ wickedness that YHWH
		dispossesses them from before you.
	}
}
{
%	\footnote{
%		take possession. Forms of this term occur more than sixty times in
%		Deuteronomy. Here it is linked specifically to the promise to the
%		patriarchs. This brings back to mind Abraham’s exchange with YHWH just
%		before their covenant, where this word occurs three times. Abraham says
%		that he is childless, so someone else will “take possession” of what is
%		his, but God assures him that his own offspring will “take possession”
%		from him (Gen 15:3–4). The numerous repetitions of the term convey the
%		central importance of the link and continuity between the episodes of
%		the patriarchs’ lives and the fate of their descendants.
%	}
}

\Verse{5}
{
	\hDtr1{לֹא בְצִדְקָתְךָ וּבְיֹשֶׁר לְבָבְךָ אַתָּה בָא לָרֶשֶׁת אֶת אַרְצָם כִּי בְּרִשְׁעַת הַגּוֹיִם הָאֵלֶּה יְהֹוָה אֱלֹהֶיךָ מוֹרִישָׁם מִפָּנֶיךָ וּלְמַעַן הָקִים אֶת הַדָּבָר אֲשֶׁר נִשְׁבַּע יְהֹוָה לַאֲבֹתֶיךָ לְאַבְרָהָם לְיִצְחָק וּלְיַעֲקֹב׃}
}
{
	\eDtr1{
		It’s not because of your virtue and your heart’s integrity
		that you’re coming to take possession of their land, but
		it is because of these nations’ wickedness that YHWH, your
		God, dispossesses them from before you, and in order to
		uphold the thing that YHWH swore to your fathers, to
		Abraham, to Isaac, and to Jacob.
	}
}
{
%	\footnote{
%		not because of your virtue. The point is made three times, in three
%		verses in a row (9:4–6): it is not because of your virtue but rather
%		because of these nations’ wickedness. Some think that the repetition is
%		just a scribe’s error, repeating a line by mistake. (Such errors are
%		known as dittography.) That is possible, but the repetition itself is
%		not sufficient reason to make that conclusion. On the contrary, the text
%		seems to me precisely to be making the point as emphatically as
%		possible. Moses notes the people’s own lack of virtue even more strongly
%		in the next verse (7) and he begins with the words “Remember—don’t
%		forget!” which is redundant as well but is certainly done on purpose for
%		emphasis. And he then goes on for the rest of the chapter listing the
%		people’s record of rebellions. Moses thus gives a powerful warning
%		against chauvinism and self-congratulation. And this also provides a
%		profound balance to the declaration that Israel was chosen to become a
%		treasured people, which came just two chapters earlier (7:6). Possession
%		of the land is the result of a promise to Israel’s ancestors. Status as
%		a treasured people depends on actions: faithfulness to the covenant.
%		Israel is not intrinsically better than anyone. What is special about
%		Israel is rather that it has been given a singular opportunity to follow
%		a path that will ultimately bring blessing to all the families of the
%		earth. Footnote 9:5. take possession. Forms of this term occur more than
%		sixty times in Deuteronomy. Here it is linked specifically to the
%		promise to the patriarchs. This brings back to mind Abraham’s exchange
%		with YHWH just before their covenant, where this word occurs three
%		times. Abraham says that he is childless, so someone else will “take
%		possession” of what is his, but God assures him that his own offspring
%		will “take possession” from him (Gen 15:3–4). The numerous repetitions
%		of the term convey the central importance of the link and continuity
%		between the episodes of the patriarchs’ lives and the fate of their
%		descendants.
%	}
}

\Verse{6}
{
	\hDtr1{וְיָדַעְתָּ כִּי לֹא בְצִדְקָתְךָ יְהֹוָה אֱלֹהֶיךָ נֹתֵן לְךָ אֶת הָאָרֶץ הַטּוֹבָה הַזֹּאת לְרִשְׁתָּהּ כִּי עַם קְשֵׁה עֹרֶף אָתָּה׃}
}
{
	\eDtr1{
		So you shall know that it is not because of your virtue
		that YHWH, your God, gives you this good land to take
		possession of it, because you’re a hardnecked people.
	}
}
{
%	\footnote{
%		take possession. Forms of this term occur more than sixty times in
%		Deuteronomy. Here it is linked specifically to the promise to the
%		patriarchs. This brings back to mind Abraham’s exchange with YHWH just
%		before their covenant, where this word occurs three times. Abraham says
%		that he is childless, so someone else will “take possession” of what is
%		his, but God assures him that his own offspring will “take possession”
%		from him (Gen 15:3–4). The numerous repetitions of the term convey the
%		central importance of the link and continuity between the episodes of
%		the patriarchs’ lives and the fate of their descendants.
%	}
}

\Verse{7}
{
	\hDtr1{זְכֹר אַל תִּשְׁכַּח אֵת אֲשֶׁר הִקְצַפְתָּ אֶת יְהֹוָה אֱלֹהֶיךָ בַּמִּדְבָּר לְמִן הַיּוֹם אֲשֶׁר יָצָאתָ מֵאֶרֶץ מִצְרַיִם עַד בֹּאֲכֶם עַד הַמָּקוֹם הַזֶּה מַמְרִים הֱיִיתֶם עִם יְהֹוָה׃}
}
{
	\eDtr1{
		Remember—don’t forget!—that you made YHWH, your God, angry
		in the wilderness. From the day that you went out from the
		land of Egypt until you came to this place, you’ve been
		rebellious toward YHWH.
	}
}
{
}

\Verse{8}
{
	\hDtr1{וּבְחֹרֵב הִקְצַפְתֶּם אֶת יְהֹוָה וַיִּתְאַנַּף יְהֹוָה בָּכֶם לְהַשְׁמִיד אֶתְכֶם׃}
}
{
	\eDtr1{
		And you made YHWH angry at Horeb, and YHWH was so incensed
		at you as to destroy you.
	}
}
{
}

\Verse{9}
{
	\hDtr1{בַּעֲלֹתִי הָהָרָה לָקַחַת לוּחֹת הָאֲבָנִים לוּחֹת הַבְּרִית אֲשֶׁר כָּרַת יְהֹוָה עִמָּכֶם וָאֵשֵׁב בָּהָר אַרְבָּעִים יוֹם וְאַרְבָּעִים לַיְלָה לֶחֶם לֹא אָכַלְתִּי וּמַיִם לֹא שָׁתִיתִי׃}
}
{
	\eDtr1{
		When I went up to the mountain to get the tablets of
		stones, the tablets of the covenant that YHWH made with
		you, and I stayed in the mountain forty days and forty
		nights, I didn’t eat bread and didn’t drink water.
	}
}
{
%	\footnote{
%		didn’t eat bread and didn’t drink water. This is the first mention of
%		this miracle that Moses personally experiences at Horeb. It is not
%		reported in Exodus. Moses tells the people about it now, in his list of
%		their rebellions, seemingly to convey to them that they rebelled in the
%		golden calf episode precisely while something awesome was going on.
%		While they were saying, “This man Moses, we don’t know what’s become of
%		him,” Moses was in fact experiencing a kind of physical transformation.
%		Another explanation: The people’s rebellions up to that point (and
%		frequently thereafter) were about food and water: the bitter water at
%		Marah (Exod 22–26), the manna (16:2–36), and the water from the rock
%		(17:1–7). And so Moses conveys to them that their worries and complaints
%		were utterly unnecessary: in God’s hands a man could even live without
%		food or water for forty days. Indeed, this follows Moses’ admonition in
%		the preceding chapter that “a human doesn’t live by bread alone” (Deut
%		8:3).
%	}
}

\Verse{10}
{
	\hDtr1{וַיִּתֵּן יְהֹוָה אֵלַי אֶת שְׁנֵי לוּחֹת הָאֲבָנִים כְּתֻבִים בְּאֶצְבַּע אֱלֹהִים וַעֲלֵיהֶם כְּכל הַדְּבָרִים אֲשֶׁר דִּבֶּר יְהֹוָה עִמָּכֶם בָּהָר מִתּוֹךְ הָאֵשׁ בְּיוֹם הַקָּהָל׃}
}
{
	\eDtr1{
		And YHWH gave me the two tablets of stones, written with
		God’s finger, and on them were all the words that YHWH had
		spoken with you at the mountain from inside the fire in
		the day of the assembly.
	}
}
{
}

\Verse{11}
{
	\hDtr1{וַיְהִי מִקֵּץ אַרְבָּעִים יוֹם וְאַרְבָּעִים לָיְלָה נָתַן יְהֹוָה אֵלַי אֶת שְׁנֵי לֻחֹת הָאֲבָנִים לֻחוֹת הַבְּרִית׃}
}
{
	\eDtr1{
		And it was: at the end of forty days and forty nights YHWH
		gave me the two tablets of stones, the tablets of the
		covenant.
	}
}
{
}

\Verse{12}
{
	\hDtr1{וַיֹּאמֶר יְהֹוָה אֵלַי קוּם רֵד מַהֵר מִזֶּה כִּי שִׁחֵת עַמְּךָ אֲשֶׁר הוֹצֵאתָ מִמִּצְרָיִם סָרוּ מַהֵר מִן הַדֶּרֶךְ אֲשֶׁר צִוִּיתִם עָשׂוּ לָהֶם מַסֵּכָה׃}
}
{
	\eDtr1{
		“And YHWH said to me, ‘Get up. Go down quickly from here,
		because your people whom you brought out from Egypt has
		become corrupt. They’ve turned quickly from the way that I
		commanded them. They’ve made themselves a molten thing.’
	}
}
{
}

\Verse{13}
{
	\hDtr1{וַיֹּאמֶר יְהֹוָה אֵלַי לֵאמֹר רָאִיתִי אֶת הָעָם הַזֶּה וְהִנֵּה עַם קְשֵׁה עֹרֶף הוּא׃}
}
{
	\eDtr1{
		And YHWH said to me, saying, ‘I’ve seen this people; and,
		here, it’s a hardnecked people.
	}
}
{
}

\Verse{14}
{
	\hDtr1{הֶרֶף מִמֶּנִּי וְאַשְׁמִידֵם וְאֶמְחֶה אֶת שְׁמָם מִתַּחַת הַשָּׁמָיִם וְאֶעֱשֶׂה אוֹתְךָ לְגוֹי עָצוּם וָרָב מִמֶּנּוּ׃}
}
{
	\eDtr1{
		Hold back from me, and I’ll destroy them and wipe out
		their name from under the skies, and I’ll make you into a
		more powerful and numerous people than they.’
	}
}
{
}

\Verse{15}
{
	\hDtr1{וָאֵפֶן וָאֵרֵד מִן הָהָר וְהָהָר בֹּעֵר בָּאֵשׁ וּשְׁנֵי לוּחֹת הַבְּרִית עַל שְׁתֵּי יָדָי׃}
}
{
	\eDtr1{
		“And I turned and went down from the mountain—and the
		mountain was burning in fire—and the two tablets of the
		covenant were on my two hands.
	}
}
{
}

\Verse{16}
{
	\hDtr1{וָאֵרֶא וְהִנֵּה חֲטָאתֶם לַיהֹוָה אֱלֹהֵיכֶם עֲשִׂיתֶם לָכֶם עֵגֶל מַסֵּכָה סַרְתֶּם מַהֵר מִן הַדֶּרֶךְ אֲשֶׁר צִוָּה יְהֹוָה אֶתְכֶם׃}
}
{
	\eDtr1{
		And I saw! And, here, you’d sinned against YHWH, your God.
		You’d made yourselves a molten calf. You’d turned quickly
		from the way that YHWH had commanded you.
	}
}
{
}

\Verse{17}
{
	\hDtr1{וָאֶתְפֹּשׂ בִּשְׁנֵי הַלֻּחֹת וָאַשְׁלִכֵם מֵעַל שְׁתֵּי יָדָי וָאֲשַׁבְּרֵם לְעֵינֵיכֶם׃}
}
{
	\eDtr1{
		And I grasped the two tablets and threw them from on my
		two hands and shattered them before your eyes.
	}
}
{
}

\Verse{18}
{
	\hDtr1{וָאֶתְנַפַּל לִפְנֵי יְהֹוָה כָּרִאשֹׁנָה אַרְבָּעִים יוֹם וְאַרְבָּעִים לַיְלָה לֶחֶם לֹא אָכַלְתִּי וּמַיִם לֹא שָׁתִיתִי עַל כּל חַטַּאתְכֶם אֲשֶׁר חֲטָאתֶם לַעֲשׂוֹת הָרַע בְּעֵינֵי יְהֹוָה לְהַכְעִיסוֹ׃}
}
{
	\eDtr1{
		“And I prostrated myself in front of YHWH like the first
		time, forty days and forty nights—I didn’t eat bread and
		didn’t drink water—over all your sin that you did, to do
		what was bad in YHWH’s eyes, to provoke Him,
	}
}
{
}

\Verse{19}
{
	\hDtr1{כִּי יָגֹרְתִּי מִפְּנֵי הָאַף וְהַחֵמָה אֲשֶׁר קָצַף יְהֹוָה עֲלֵיכֶם לְהַשְׁמִיד אֶתְכֶם וַיִּשְׁמַע יְהֹוָה אֵלַי גַּם בַּפַּעַם הַהִוא׃}
}
{
	\eDtr1{
		because I was dreading on account of the anger and the
		fury, that YHWH was so angry at you as to destroy you. But
		YHWH listened to me that time as well.
	}
}
{
}

\Verse{20}
{
	\hDtr1{וּבְאַהֲרֹן הִתְאַנַּף יְהֹוָה מְאֹד לְהַשְׁמִידוֹ וָאֶתְפַּלֵּל גַּם בְּעַד אַהֲרֹן בָּעֵת הַהִוא׃}
}
{
	\eDtr1{
		And YHWH was very incensed at Aaron so as to destroy him,
		and I prayed for Aaron at that time as well.
	}
}
{
%	\footnote{
%		incensed at Aaron. The text of the golden calf event (Exodus 32) never
%		says that God was specifically angry at Aaron. We might understand the
%		fact that Moses only makes it known now for the first time as indicating
%		that he would not mention it until after Aaron’s death. Thus he spares
%		Aaron the humiliation. And he has thus deferred any doubts that this
%		might have raised about the status of the high priest until now, when
%		they can be considered in the light of the man’s entire life and career
%		as high priest.
%	}
}

\Verse{21}
{
	\hDtr1{וְאֶת חַטַּאתְכֶם אֲשֶׁר עֲשִׂיתֶם אֶת הָעֵגֶל לָקַחְתִּי וָאֶשְׂרֹף אֹתוֹ בָּאֵשׁ וָאֶכֹּת אֹתוֹ טָחוֹן הֵיטֵב עַד אֲשֶׁר דַּק לְעָפָר וָאַשְׁלִךְ אֶת עֲפָרוֹ אֶל הַנַּחַל הַיֹּרֵד מִן הָהָר׃}
}
{
	\eDtr1{
		And your sin that you made, the calf: I took it and burned
		it in fire and crushed it, grinding it well until it was
		thin as dust, and I threw its dust into the wadi that
		comes down from the mountain.
	}
}
{
%	\footnote{
%		burned it … thin as dust. The parallels between Aaron’s golden calf and
%		those of Jeroboam I, king of Israel, are unmistakable. Jeroboam erects
%		golden calves at Dan and Beth-El (1 Kings 12:26–30). On that occasion he
%		says, “Here are your gods, Israel, which brought you out from the land
%		of Egypt”—which are the same words that the people say at the golden
%		calf in Exodus. Jeroboam’s sons are Nadab and Abiyah; Aaron’s sons are
%		Nadab and Abihu. And King Josiah destroys the high place where
%		Jeroboam’s golden calf stood at Beth-El, and, like Moses, “he burned it
%		thin as dust” (2 Kings 23:15). Also like Moses, Josiah casts the dust of
%		pagan altars into a wadi (23:6,12). A literary-historical analysis of
%		this parallel appears in my Who Wrote the Bible? (pp. 111–116). In the
%		terms of the full Tanak as it now stands, the parallel would suggest
%		that history repeats itself. This is also one of many parallels between
%		Josiah and Moses that serve to single out Josiah as the best of the
%		kings of Israel, the one who comes closest to the standards that Moses
%		taught by commandment and example. Thus the Torah says at the end of
%		Moses’ life: “And a prophet did not rise again in Israel like Moses”
%		(Deut 34:10); and the Tanak says at the end of Josiah’s life: “And after
%		him one did not rise like him” (2 Kings 23:25)—and these are the only
%		two occurrences of the expression “did not rise like him” in the Bible.
%		Again the embroidery of connections through the Tanak is apparent.
%	}
}

\Verse{22}
{
	\hDtr1{וּבְתַבְעֵרָה וּבְמַסָּה וּבְקִבְרֹת הַתַּאֲוָה מַקְצִפִים הֱיִיתֶם אֶת יְהֹוָה׃}
}
{
	\eDtr1{
		“And at Taberah and at Massah and at Kibroth Hattaavah you
		were making YHWH angry.
	}
}
{
}

\Verse{23}
{
	\hDtr1{וּבִשְׁלֹחַ יְהֹוָה אֶתְכֶם מִקָּדֵשׁ בַּרְנֵעַ לֵאמֹר עֲלוּ וּרְשׁוּ אֶת הָאָרֶץ אֲשֶׁר נָתַתִּי לָכֶם וַתַּמְרוּ אֶת פִּי יְהֹוָה אֱלֹהֵיכֶם וְלֹא הֶאֱמַנְתֶּם לוֹ וְלֹא שְׁמַעְתֶּם בְּקֹלוֹ׃}
}
{
	\eDtr1{
		And when YHWH sent you from Kadesh-barnea, saying, ‘Go up
		and take possession of the land that I’ve given you,’ then
		you rebelled at the word of YHWH, your God, and you didn’t
		trust Him and didn’t listen to His voice.
	}
}
{
}

\Verse{24}
{
	\hDtr1{מַמְרִים הֱיִיתֶם עִם יְהֹוָה מִיּוֹם דַּעְתִּי אֶתְכֶם׃}
}
{
	\eDtr1{You’ve been rebelling toward YHWH from the day I knew you.}
}
{
}

\Verse{25}
{
	\hDtr1{וָאֶתְנַפַּל לִפְנֵי יְהֹוָה אֵת אַרְבָּעִים הַיּוֹם וְאֶת אַרְבָּעִים הַלַּיְלָה אֲשֶׁר הִתְנַפָּלְתִּי כִּי אָמַר יְהֹוָה לְהַשְׁמִיד אֶתְכֶם׃}
}
{
	\eDtr1{
		“So I prostrated myself in front of YHWH for the forty
		days and forty nights that I had fallen down because YHWH
		had said He would destroy you,
	}
}
{
}

\Verse{26}
{
	\hDtr1{וָאֶתְפַּלֵּל אֶל יְהֹוָה וָאֹמַר אֲדֹנָי יֱהֹוִה אַל תַּשְׁחֵת עַמְּךָ וְנַחֲלָתְךָ אֲשֶׁר פָּדִיתָ בְּגדְלֶךָ אֲשֶׁר הוֹצֵאתָ מִמִּצְרַיִם בְּיָד חֲזָקָה׃}
}
{
	\eDtr1{
		and I prayed to YHWH and said, ‘My Lord YHWH, don’t
		destroy your people and your legacy whom you redeemed by
		your greatness, whom you brought out from Egypt with a
		strong hand.
	}
}
{
%	\footnote{
%		destroy. The word in Hebrew, ta[image][image][image]t, has the same root
%		meaning as the word for what the people have done: become corrupt,
%		hi[image][image]ît (9:12). The same play on this root is used in the
%		flood story (Gen 6:12–13). There and here it is not just wordplay for
%		its own sake. It conveys the principle that punishment is meant to be in
%		proportion to the crime. The repetition of the pun here in Deuteronomy
%		conveys that what was true for all humankind in the early generations of
%		creation is still true in relations between God and humans in the
%		Torah’s final book.
%	}
}

\Verse{27}
{
	\hDtr1{זְכֹר לַעֲבָדֶיךָ לְאַבְרָהָם לְיִצְחָק וּלְיַעֲקֹב אַל תֵּפֶן אֶל קְשִׁי הָעָם הַזֶּה וְאֶל רִשְׁעוֹ וְאֶל חַטָּאתוֹ׃}
}
{
	\eDtr1{
		Remember your servants Abraham, Isaac, and Jacob. Don’t
		look to this people’s hardness and to its wickedness and
		to its sin,
	}
}
{
}

\Verse{28}
{
	\hDtr1{פֶּן יֹאמְרוּ הָאָרֶץ אֲשֶׁר הוֹצֵאתָנוּ מִשָּׁם מִבְּלִי יְכֹלֶת יְהֹוָה לַהֲבִיאָם אֶל הָאָרֶץ אֲשֶׁר דִּבֶּר לָהֶם וּמִשִּׂנְאָתוֹ אוֹתָם הוֹצִיאָם לַהֲמִתָם בַּמִּדְבָּר׃}
}
{
	\eDtr1{
		or else the land from which you brought us out will say:
		Because YHWH wasn’t able to bring them to the land of
		which He spoke to them, and because He hated them, He
		brought them out to kill them in the wilderness!
	}
}
{
}

\Verse{29}
{
	\hDtr1{וְהֵם עַמְּךָ וְנַחֲלָתֶךָ אֲשֶׁר הוֹצֵאתָ בְּכֹחֲךָ הַגָּדֹל וּבִזְרֹעֲךָ הַנְּטוּיָה׃}
}
{
	\eDtr1{
		And they’re your people and your legacy, whom you brought
		out by your great power and by your outstretched arm.’
	}
}
{
}
